% ===========================================================================
% Appendix: Formal Specification for Mechanized Verification
% ===========================================================================
\section{Formal Specification for Mechanized Verification}
\label{app:coq}

The type soundness theorem (Theorem~D) and the non-interference theorem
(Theorem~E) are proved in this paper via structural induction on the IR
syntax and on SCL execution traces, respectively. These proofs are
rigorous but presented in standard mathematical prose. For readers
interested in machine-checked verification, we provide the formal
specification of the key lemmas in a notation compatible with Coq
and Isabelle/HOL, together with a mechanization of the core SVNN type
system and the repaired Galois connection
(\texttt{code/coq/neuroplc\_core.v}, compiled with Coq~9.1, 2026-08-05):
the typing judgment, the interval domain, the soundness direction of
the Galois pair (any certificate value lies in the concrete envelope,
whose $M_2$-term bounds the LUT error), and an abstract type-soundness
statement. A mechanization of the balanced-allocation closure of
Theorem~\ref{thm:optimal_lut} (\texttt{code/coq/optimal\_allocation.v})
reduces the minimax optimality to a convexity argument (the balanced
point equalizes the two errors; any other split makes at least one
error strictly larger).

\subsection{IR Syntax (Inductive Type)}

\begin{verbatim}
Inductive ir_expr : Type :=
  | IRInput  : var -> ir_expr
  | IRMatMul : matrix -> ir_expr -> ir_expr
  | IRBsplineLUT : lut_table -> grid -> ir_expr -> ir_expr
  | IRStandardAct : activation -> ir_expr -> ir_expr
  | IRAdd    : ir_expr -> ir_expr -> ir_expr
  | IRSoftmax : ir_expr -> ir_expr
  | IRArgmax : ir_expr -> ir_expr.
\end{verbatim}

\subsection{Type System (Inductive Typing Judgment)}

\begin{verbatim}
Inductive ir_type : Type :=
  | LinearT   : ir_type
  | ElemwiseT : ir_type
  | SVNNT     : ir_type.

Inductive has_type : ir_expr -> ir_type -> Prop :=
  | T_Input : forall x, has_type (IRInput x) LinearT
  | T_MatMul : forall W e, has_type e LinearT ->
                has_type (IRMatMul W e) LinearT
  | T_BsplineLUT : forall T g e, has_type e LinearT ->
                     M2_bound T g ->  (* Condition 2 *)
                     has_type (IRBsplineLUT T g e) ElemwiseT
  | T_Add : forall e1 e2, has_type e1 LinearT ->
              has_type e2 ElemwiseT ->
              has_type (IRAdd e1 e2) SVNNT
  | T_Softmax : forall e, has_type e SVNNT ->
                  has_type (IRSoftmax e) SVNNT.
\end{verbatim}

\subsection{Type Soundness Statement}

\begin{verbatim}
Theorem type_soundness :
  forall (e : ir_expr) (rho : env),
    has_type e SVNNT ->
    exists (eps : R),
      eps > 0 /\
      forall (x : input_vector),
        in_domain x ->
        Rabs (denote_real e rho x - denote_scl (compile e) rho x)
          <= eps.
\end{verbatim}

\textit{Proof sketch (structural induction on \texttt{has\_type}):}
Base cases (\texttt{T\_Input}, \texttt{T\_MatMul}, \texttt{T\_BsplineLUT})
use the de~Boor error bound (Lemma~\texttt{lut\_error\_bound}) and DA
propagation (Lemma~\texttt{da\_linear\_soundness}). Inductive cases
(\texttt{T\_Add}, \texttt{T\_Softmax}) compose error bounds using the
Lipschitz property (Lemma~\texttt{add\_lipschitz},
Lemma~\texttt{softmax\_\allowbreak{}contraction}).

\subsection{Non-Interference Statement}

\begin{verbatim}
Theorem non_interference :
  forall (P : scl_program) (N : svnn_network) (phi : safety_property),
    satisfies P phi ->
    let F := compile_N N in
    let P_combined := compose P F in
    satisfies P_combined phi.
\end{verbatim}

\textit{Proof sketch (induction on SCL trace length):}
Base case ($n=0$): initial states satisfy $\varphi$ by assumption.
Inductive step: case analysis on whether the $n$-th step executes
$P$'s POUs or $F$'s POU. Proposition~NC.1 (memory isolation) ensures
$\text{State}_P$ is invariant when $F$ executes. Proposition~NC.2
(termination) bounds $F$'s execution within one scan cycle.
Proposition~NC.3 (numerical safety) guarantees $F$'s outputs are
finite. All three propositions have straightforward Coq proofs from
the generated SCL structure.

\subsection{DA Galois Connection (Semantic Domain)}

\begin{verbatim}
Record doubleton := mk_doubleton { center : R ; radius : R }.
(* Weighted first-order-plus-curvature envelope; the right adjoint
   generated by alpha (radius includes the first-order term). *)
Definition gamma (d : doubleton) (f : R -> R) (x0 r : R) : Prop :=
  Rabs (f x0 - center d) + Rabs (f' x0) * r + M2 f * r * r / 2
    <= radius d.
Definition alpha (f : R -> R) (x0 r : R) : doubleton :=
  mk_doubleton (f x0) (Rabs (f' x0) * r + M2 f * r * r / 2).

Lemma galois_soundness :
  forall f x0 r, gamma (alpha f x0 r) f x0 r.
(* Tightness: the linear witness f(x)=c+(R/r)(x-x0) attains the
   radius, so alpha o gamma is tight (infimum = d). The insertion
   form alpha o gamma = id does not hold (constant functions). *)
Lemma galois_tightness :
  forall c R x0 r (rpos : 0 < r),
    exists f, gamma (mk_doubleton c R) f x0 r
      /\ alpha f x0 r = mk_doubleton c R.
Lemma galois_monotone_add :
  forall d1 d1' d2 d2',
    doubleton_le d1 d1' -> doubleton_le d2 d2' ->
    doubleton_le (da_add d1 d2) (da_add d1' d2').
\end{verbatim}

\vspace{4pt}
\noindent\textbf{Mechanization status.}
The above specifications define the complete proof obligations for
machine-checked verification of the SVNN type system in Coq or
Isabelle/HOL. The proofs---structural induction on the IR syntax
(Theorem~D), induction on SCL trace length (Theorem~E), and
algebraic verification of the Galois connection properties
(Theorem~B)---follow the mathematical proofs presented in
{\S}\ref{sec:ir-semantics}, {\S}\ref{sec:noninterference}, and
{\S}\ref{sec:galois} respectively. Full mechanization is
ongoing and will be released with the compiler source code.
